%% Web Appendix
%% "Does AI Reduce the Reputation Premium in Knowledge Diffusion?
%%  A Conceptual Replication Using Chinese Journal Articles"
%% Marketing Letters -- Replication Corner

\documentclass[12pt,a4paper]{article}
\usepackage[top=2.54cm,bottom=2.54cm,left=2.54cm,right=2.54cm]{geometry}
\usepackage{setspace}
\usepackage{amsmath,amssymb}
\usepackage{booktabs}
\usepackage{threeparttable}
\usepackage{array}
\usepackage{longtable}
\usepackage[authoryear,sort]{natbib}
\usepackage[hidelinks]{hyperref}
\usepackage{microtype}
\usepackage{times}
\usepackage[utf8]{inputenc}
\usepackage[T1]{fontenc}
\usepackage{caption}
\captionsetup{labelfont=bf,font=small}
\newcommand{\doi}[1]{\href{https://doi.org/#1}{doi:#1}}

\onehalfspacing

\title{\textbf{Web Appendix}\\[0.4em]
\normalsize Does AI Reduce the Reputation Premium in Knowledge Diffusion?\\
A Conceptual Replication Using Chinese Journal Articles\\[0.4em]
\textit{Marketing Letters} --- Replication Corner}
\author{[Blinded for review]}
\date{}

\begin{document}
\maketitle
\tableofcontents
\newpage

%% ============================================================
\section*{Appendix A: Data Construction and Variable Coding}
\addcontentsline{toc}{section}{Appendix A: Data Construction and Variable Coding}
%% ============================================================

\subsection*{A1. Sample Construction}
\addcontentsline{toc}{subsection}{A1. Sample Construction}

\subsubsection*{CNKI Search Procedure}

We searched CNKI (China National Knowledge Infrastructure,
\url{https://www.cnki.net}) in June 2026 for Chinese-language journal
articles using the advanced search interface.
The search covered journals classified in CNKI's economics, management,
and business administration categories.
Keywords applied (in separate searches combined by OR logic):
人工智能 (AI), 生成式AI (generative AI), 大语言模型 (LLM), 数字化转型
(digital transformation), 公司治理 (corporate governance),
供应链管理 (supply chain management), ESG, 绿色发展 (green development),
营销 (marketing), 技术创新 (technological innovation),
企业管理 (enterprise management), 数字经济 (digital economy).
Publication date range: 2023-01-01 to 2026-06-05 (data collection date).
Document types: journal articles only (editorials, book reviews,
conference proceedings excluded).

\subsubsection*{Sample Flow}

\begin{table}[htbp]
\centering
\caption{Sample Attrition Flow}
\label{tab:flow}
\begin{tabular}{lrr}
\toprule
Step & $N$ & Reason for exclusion \\
\midrule
Initial CNKI records retrieved & 856 & \\
Records outside 2023--2026 window & $-$5 & Publication date outside scope \\
Records with missing or partial abstracts & $-$6 & $<$20 characters in abstract field \\
Records lacking formal publication ID & $-$2 & Online-first without issue/page data \\
\midrule
\textbf{Final analytic sample} & \textbf{843} & \\
\bottomrule
\end{tabular}
\end{table}

\subsubsection*{Year and Topic Distribution}

Table~\ref{tab:yeardist} shows the analytic sample by publication year.
The 2025--2026 concentration reflects both the focus on recent AI-era
publications and the exponential growth of CNKI-indexed articles in
this period.

\begin{table}[htbp]
\centering
\caption{Analytic Sample by Publication Year}
\label{tab:yeardist}
\begin{tabular}{lrrr}
\toprule
Year & $N$ & \% & Cumulative \% \\
\midrule
2023 & 100 & 11.9 & 11.9 \\
2024 & 151 & 17.9 & 29.8 \\
2025 & 395 & 46.9 & 76.6 \\
2026 & 197 & 23.4 & 100.0 \\
\midrule
Total & 843 & 100.0 & \\
\bottomrule
\end{tabular}
\end{table}

The keyword strategy was designed to retrieve articles relevant to
business and management knowledge diffusion; 88\% of articles were
classified as AI/digital transformation-related (see topic
classification in Appendix~A4).
A marketing and management sub-sample robustness check is reported
in Appendix~B7.

%% ============================================================
\subsection*{A2. AI-Style Score: Construction and Validation}
\addcontentsline{toc}{subsection}{A2. AI-Style Score}
%% ============================================================

\subsubsection*{Dictionary Construction}

Two bilingual coders (Chinese-English) independently reviewed a
development sample of 120 abstracts from 2023--2024, half of which
were identified as likely AI-assisted (based on a preliminary scan
for formulaic structural features) and half not.
Using the AI writing-style taxonomy of Liang et al.\ (2024, \textit{NEJM AI})
as a starting framework, coders independently proposed candidate markers
characteristic of AI-assisted academic writing in Chinese social science.
Candidate markers were operationalized as exact-string matches on the
abstract text.
After resolving discrepancies by discussion, a final dictionary of
23 markers was agreed upon.
Inter-coder agreement on the development sample was $\kappa = 0.79$,
exceeding the conventional 0.70 acceptability threshold.

\subsubsection*{Scoring Procedure}

For each target abstract, we assessed the binary presence (1/0) of
each of the 23 markers via exact string matching.
The \textit{ai\_style\_score} equals the proportion of markers matched:
\[
  \textit{ai\_style\_score}_i
    = \frac{1}{23}\sum_{j=1}^{23} \mathbf{1}[\text{marker}_j \in \text{abstract}_i]
  \in [0,1].
\]
This fully deterministic scoring procedure requires no external API
and is entirely reproducible from the published marker list.

\subsubsection*{Limitations of the Measure}

CNKI truncates the displayed abstract field at approximately 203 characters,
which means abstracts are right-truncated: phrases that would appear in
longer sections of an abstract are systematically missing.
This truncation likely suppresses scores for markers that appear in
later structural sections (e.g., practical implications) more than for
markers that appear early (e.g., ``results show'').
As Table~\ref{tab:markers} shows, several markers (理论贡献,
实践启示, 数据表明, 研究发现如下, 系统性地) have zero frequency in
the analytic sample, plausibly reflecting truncation rather than
genuine absence.
Future work should validate the score against full-text abstracts or
alternative AI-detection approaches.

\subsubsection*{Full Marker List with Frequencies}

\begin{table}[htbp]
\centering
\caption{AI Writing-Style Marker Dictionary: Chinese Markers,
         English Translations, and Frequency in Analytic Sample
         ($N = 843$)}
\label{tab:markers}
\begin{threeparttable}
\small
\begin{tabular}{llrr}
\toprule
\# & Chinese marker & English translation & $n$ (\%) \\
\midrule
1  & 研究发现  & results/findings indicate & 338 (40.1) \\
2  & 结果表明  & results show              & 108 (12.8) \\
3  & 异质性分析 & heterogeneity analysis   &  80 (9.5)  \\
4  & 稳健性检验 & robustness checks        &  77 (9.1)  \\
5  & 机制分析  & mechanism analysis        &  72 (8.5)  \\
6  & 研究表明  & research shows            &  42 (5.0)  \\
7  & 面板数据  & panel data                &  38 (4.5)  \\
8  & 双重差分  & difference-in-differences &  36 (4.3)  \\
9  & 本研究    & this study                &  35 (4.2)  \\
10 & 固定效应  & fixed effects             &  27 (3.2)  \\
11 & 实证研究  & empirical study           &  22 (2.6)  \\
12 & 中介效应  & mediation effect          &  16 (1.9)  \\
13 & 调节效应  & moderation effect         &  16 (1.9)  \\
14 & 进一步分析 & further analysis         &  16 (1.9)  \\
15 & 创新性地  & innovatively              &   5 (0.6)  \\
16 & 验证了    & validates / confirms      &   5 (0.6)  \\
17 & 研究结论  & research conclusions      &   7 (0.8)  \\
18 & 基准回归  & baseline regression       &   6 (0.7)  \\
19 & 理论贡献  & theoretical contribution  &   0 (0.0)  \\
20 & 实践启示  & practical implications    &   0 (0.0)  \\
21 & 数据表明  & data show                 &   0 (0.0)  \\
22 & 研究发现如下 & findings are as follows &  0 (0.0)  \\
23 & 系统性地  & systematically            &   0 (0.0)  \\
\midrule
\multicolumn{2}{l}{\textit{Total marker occurrences}} & 931 & --- \\
\bottomrule
\end{tabular}
\begin{tablenotes}
\footnotesize
\item \textit{Note.} Frequencies reflect binary presence in each abstract;
  multiple matches within one abstract count as one.
  Nine markers with zero or near-zero frequency likely reflect CNKI
  abstract truncation at $\approx$203 characters rather than genuine
  absence of these patterns from AI-assisted writing.
  All matching is exact string search; no stemming or paraphrase matching.
\end{tablenotes}
\end{threeparttable}
\end{table}

\subsubsection*{AI-Style Score Distribution and Temporal Trend}

The score is right-skewed: 40.6\% of articles score exactly zero
(no markers present), while the maximum is 0.261
(6 of 23 markers). The mean is 0.049 (SD\,=\,0.050, median\,=\,0.044).

Table~\ref{tab:aitrend} shows mean AI-style scores by publication year.
The increasing trend from 2023 to 2026 is consistent with the
accelerating adoption of LLM-assisted writing following the release
of ChatGPT in late 2022.

\begin{table}[htbp]
\centering
\caption{Mean AI-Style Score by Publication Year}
\label{tab:aitrend}
\begin{tabular}{lrrr}
\toprule
Year & Mean & Median & SD \\
\midrule
2023 & 0.031 & 0.022 & 0.041 \\
2024 & 0.041 & 0.044 & 0.047 \\
2025 & 0.050 & 0.044 & 0.049 \\
2026 & 0.062 & 0.044 & 0.056 \\
\bottomrule
\end{tabular}
\end{table}

%% ============================================================
\subsection*{A3. Institutional Prestige Coding and First-Author Sensitivity}
\addcontentsline{toc}{subsection}{A3. Prestige Coding and Sensitivity}
%% ============================================================

\subsubsection*{Prestige Classification}

Table~\ref{tab:prestgedist} shows the prestige tier distribution.

\begin{table}[htbp]
\centering
\caption{Institutional Prestige Tier Distribution ($N = 843$)}
\label{tab:prestgedist}
\begin{tabular}{llrr}
\toprule
Tier & Classification & $N$ & \% \\
\midrule
1 & No nationally ranked institution & 280 & 33.2 \\
2 & 211 Project / Double First-Class & 191 & 22.7 \\
3 & 985 Project institution          & 246 & 29.2 \\
4 & C9 elite university              & 126 & 14.9 \\
\midrule
\multicolumn{2}{l}{Prestige\_High ($\geq 2$)} & 563 & 66.8 \\
\bottomrule
\end{tabular}
\end{table}

The maximum-tier rule assigns each article the highest prestige tier
across all affiliated co-authors.
Accordingly, 51.7\% of articles have a high-prestige \textit{first} author
(versus 66.8\% when using maximum tier), reflecting that some articles
gain high-prestige classification only through non-first-author affiliations.

\subsubsection*{First-Author Prestige Sensitivity}

Table~\ref{tab:firstauthor} replicates the main OLS specifications
using first-author institutional prestige rather than maximum-tier
prestige.
Notably, with first-author coding the prestige main effect
on downloads reaches conventional significance
($\hat\beta = 0.157$, $p = .017$), suggesting that the maximum-tier
rule attenuates the prestige effect by assigning high-prestige status
to articles whose first author is not from an elite institution.
Importantly, the Prestige\,$\times$\,AI-style interaction remains
non-significant in both specifications
($\hat\beta_3 = -0.024$, $p = .699$ for downloads;
$\hat\beta_3 = +0.086$, $p = .115$ for citations),
consistent with the null interaction finding in the main text.
These results confirm that the null interaction is robust to the
choice of prestige aggregation rule, even when the prestige main
effect itself is significant under first-author coding.

\begin{table}[htbp]
\centering
\caption{Sensitivity Analysis: First-Author Institutional Prestige
         ($N = 843$, OLS, HC3 Robust SEs)}
\label{tab:firstauthor}
\begin{threeparttable}
\small
\begin{tabular}{lcc}
\toprule
 & (1) Downloads & (2) Citations \\
\midrule
Prestige\_High\_first
  & $0.157^{*}$ & $0.063$ \\
  & (0.065)     & (0.061) \\[4pt]
ai\_z
  & $0.215^{***}$ & $0.060$ \\
  & (0.047)       & (0.043) \\[4pt]
\textbf{Prestige\_first\,$\times$\,ai\_z}
  & $\mathbf{-0.024}$ & $\mathbf{+0.086}$ \\
  & \textbf{(0.062)}  & \textbf{(0.054)}  \\[4pt]
\midrule
Controls, Year FEs & Yes & Yes \\
$R^2$              & .370 & .666 \\
$N$                & 843  & 843  \\
\bottomrule
\end{tabular}
\begin{tablenotes}
\footnotesize
\item \textit{Note.} Prestige\_High\_first\,=\,1 if the first-listed author is
  affiliated with a 985/C9/211/DFC-ranked institution (51.7\% of sample).
  Controls include $\log(\text{Age})$, $\log(\text{Abs.\ length})$,
  and year fixed effects.
  $^{*}p<.05$;\ $^{***}p<.001$.
\end{tablenotes}
\end{threeparttable}
\end{table}

%% ============================================================
\subsection*{A4. Variable Dictionary and Additional Descriptives}
\addcontentsline{toc}{subsection}{A4. Variable Dictionary}
%% ============================================================

\begin{table}[htbp]
\centering
\caption{Variable Definitions and Summary Statistics ($N = 843$)}
\label{tab:vardict}
\begin{threeparttable}
\small
\begin{tabular}{lp{5.5cm}rr}
\toprule
Variable & Definition & Mean & SD \\
\midrule
$\log(\text{Downloads}+1)$ & Log of CNKI download count plus one & 7.52 & 1.19 \\
$\log(\text{Citations}+1)$ & Log of CNKI citation count plus one & 1.81 & 1.52 \\
Prestige\_High & 1 if max tier $\geq 2$ (985/C9/211/DFC) & 0.668 & 0.471 \\
prestige\_score\_z & Standardized four-tier prestige score & 0.000 & 1.000 \\
ai\_score & Proportion of 23 AI markers present & 0.049 & 0.050 \\
ai\_z & Standardized ai\_score & 0.000 & 1.000 \\
$\log(\text{Age})$ & Log of days from publication to 2026-06-05 & 5.69 & 0.94 \\
$\log(\text{Abs.\ length})$ & Log of abstract character count & 5.32 & 0.033 \\
\bottomrule
\end{tabular}
\begin{tablenotes}
\footnotesize
\item \textit{Note.} The small SD for $\log(\text{Abs.\ length})$ reflects
  CNKI truncation at $\approx$203 characters; see main text and
  Table~\ref{tab:descriptives} in the manuscript.
\end{tablenotes}
\end{threeparttable}
\end{table}

\newpage
%% ============================================================
\section*{Appendix B: Additional Analyses}
\addcontentsline{toc}{section}{Appendix B: Additional Analyses}
%% ============================================================

\subsection*{B1. Full Main Regression Tables}
\addcontentsline{toc}{subsection}{B1. Full Main Regression Tables}

Tables~\ref{tab:b1_downloads} and~\ref{tab:b1_citations} reproduce
the main regression results with all coefficients including intercept,
year fixed effects, and control variables.

\begin{table}[htbp]
\centering
\caption{Full Regression Output: Download Outcome ($N = 843$)}
\label{tab:b1_downloads}
\begin{threeparttable}
\small
\begin{tabular}{lcc}
\toprule
 & (1) Binary prestige & (3) Continuous prestige \\
\midrule
Intercept
  & $-6.506$ & $-6.623$ \\
  & (8.564)  & (8.599)  \\[3pt]
Prestige\_High
  & $0.113$ & \\
  & (0.069) & \\[3pt]
prestige\_score\_z
  & & $0.031$ \\
  & & (0.035) \\[3pt]
ai\_z
  & $0.216^{***}$ & $0.187^{***}$ \\
  & (0.048)       & (0.032)       \\[3pt]
\textbf{Prestige\,$\times$\,ai\_z}
  & $\mathbf{-0.044}$ & $\mathbf{-0.013}$ \\
  & \textbf{(0.062)}  & \textbf{(0.032)}  \\[3pt]
$\log(\text{Age})$
  & $0.697^{***}$ & $0.702^{***}$ \\
  & (0.090)       & (0.090)       \\[3pt]
$\log(\text{Abs.\ length})$
  & $1.825$ & $1.876$ \\
  & (1.601) & (1.629) \\[3pt]
Year = 2024
  & $0.452^{***}$ & $0.449^{***}$ \\
  & (0.118)       & (0.118)       \\[3pt]
Year = 2025
  & $0.391^{**}$ & $0.384^{**}$ \\
  & (0.143)      & (0.143)      \\[3pt]
Year = 2026
  & $0.077$ & $0.069$ \\
  & (0.248) & (0.249) \\[3pt]
\midrule
$R^2$       & .368 & .366 \\
Adj.\ $R^2$ & .362 & .360 \\
$N$         & 843  & 843  \\
\bottomrule
\end{tabular}
\begin{tablenotes}
\footnotesize
\item \textit{Note.} OLS with HC3 heteroskedasticity-robust SEs in parentheses.
  $^{**}p<.01$;\ $^{***}p<.001$.
  Baseline year = 2023.
\end{tablenotes}
\end{threeparttable}
\end{table}

\begin{table}[htbp]
\centering
\caption{Full Regression Output: Citation Outcome ($N = 843$)}
\label{tab:b1_citations}
\begin{threeparttable}
\small
\begin{tabular}{lcc}
\toprule
 & (2) Binary prestige & (4) Continuous prestige \\
\midrule
Intercept
  & $-9.939$ & $-9.906$ \\
  & (11.886) & (11.917) \\[3pt]
Prestige\_High
  & $0.046$ & \\
  & (0.065) & \\[3pt]
prestige\_score\_z
  & & $0.016$ \\
  & & (0.031) \\[3pt]
ai\_z
  & $0.059$ & $0.065^{*}$ \\
  & (0.044) & (0.028)    \\[3pt]
\textbf{Prestige\,$\times$\,ai\_z}
  & $\mathbf{+0.010}$ & $\mathbf{+0.002}$ \\
  & \textbf{(0.056)}  & \textbf{(0.028)}  \\[3pt]
$\log(\text{Age})$
  & $0.900^{***}$ & $0.901^{***}$ \\
  & (0.093)       & (0.093)       \\[3pt]
$\log(\text{Abs.\ length})$
  & $1.395$ & $1.404$ \\
  & (2.224) & (2.228) \\[3pt]
Year = 2024
  & $-0.178$ & $-0.177$ \\
  & (0.139)  & (0.139)  \\[3pt]
Year = 2025
  & $-1.057^{***}$ & $-1.054^{***}$ \\
  & (0.163)        & (0.163)        \\[3pt]
Year = 2026
  & $-1.251^{***}$ & $-1.249^{***}$ \\
  & (0.251)        & (0.251)        \\[3pt]
\midrule
$R^2$       & .663 & .663 \\
Adj.\ $R^2$ & .660 & .660 \\
$N$         & 843  & 843  \\
\bottomrule
\end{tabular}
\begin{tablenotes}
\footnotesize
\item \textit{Note.} OLS with HC3 SEs. The large negative year fixed effects
  for 2025--2026 reflect citation left-censoring: recently published
  articles have had little time to accumulate citations.
  $^{*}p<.05$;\ $^{***}p<.001$.
\end{tablenotes}
\end{threeparttable}
\end{table}

%% ============================================================
\subsection*{B2. Full Robustness Regression Tables}
\addcontentsline{toc}{subsection}{B2. Full Robustness Regression Tables}

Table~\ref{tab:b2_robust} provides key coefficients for all six
robustness specifications reported in Table~3 of the main text.
All six specifications include $\log(\text{Age})$,
$\log(\text{Abs.\ length})$, and year fixed effects.

\begin{table}[htbp]
\centering
\caption{Robustness Checks: Full Key Coefficients}
\label{tab:b2_robust}
\begin{threeparttable}
\small
\setlength{\tabcolsep}{5pt}
\begin{tabular}{lp{3.5cm}cccc}
\toprule
Spec. & Description & Prestige & AI-Style & \textbf{Prestige\,$\times$\,AI} & $R^2$ \\
\midrule
1 & Baseline: binary prestige, downloads ($n=843$)
  & $0.113$ & $0.216^{***}$ & $\mathbf{-0.044}$ & .368 \\
  & & (0.069) & (0.048) & (0.062) & \\[3pt]
2 & Binary AI median-split, downloads ($n=843$)
  & $0.159^{\dagger}$ & $0.383^{***}$ & $\mathbf{-0.149}$ & .356 \\
  & & (0.089) & (0.108) & (0.139) & \\[3pt]
3 & Exclude 2026, downloads ($n=646$)
  & $0.097$ & $0.174^{***}$ & $\mathbf{-0.001}$ & .137 \\
  & & (0.073) & (0.053) & (0.068) & \\[3pt]
4 & Trim top 5\% downloads, downloads ($n=800$)
  & $0.074$ & $0.213^{***}$ & $\mathbf{-0.042}$ & .372 \\
  & & (0.065) & (0.046) & (0.059) & \\[3pt]
5 & Baseline: binary prestige, citations ($n=843$)
  & $0.046$ & $0.059$ & $\mathbf{+0.010}$ & .663 \\
  & & (0.065) & (0.044) & (0.056) & \\[3pt]
6 & Continuous prestige score, downloads ($n=843$)
  & $0.031$ & $0.187^{***}$ & $\mathbf{-0.013}$ & .366 \\
  & & (0.035) & (0.032) & (0.032) & \\[3pt]
\bottomrule
\end{tabular}
\begin{tablenotes}
\footnotesize
\item \textit{Note.} OLS with HC3 robust SEs in parentheses.
  The Prestige\,$\times$\,AI interaction (bold) is not significant
  in any specification.
  $^{\dagger}p<.10$;\ $^{***}p<.001$.
\end{tablenotes}
\end{threeparttable}
\end{table}

%% ============================================================
\subsection*{B5. Negative-Binomial Models on Raw Counts}
\addcontentsline{toc}{subsection}{B5. Negative-Binomial Models}

Table~\ref{tab:nb} presents negative-binomial regressions on raw
download and citation counts as an alternative to the log(+1)-OLS
specification in the main analysis.
Negative-binomial regression is appropriate for count outcomes with
overdispersion.
Coefficients are log-incidence-rate ratios (IRR on the exponentiated
scale in parentheses).

\begin{table}[htbp]
\centering
\caption{Negative-Binomial Models on Raw Counts ($N = 843$)}
\label{tab:nb}
\begin{threeparttable}
\small
\begin{tabular}{lcc}
\toprule
 & (1) Downloads & (2) Citations \\
\midrule
Prestige\_High
  & $0.158^{*}$   & $0.076$ \\
  & (IRR\,=\,1.171) & (IRR\,=\,1.079) \\[3pt]
ai\_z
  & $0.146^{**}$   & $0.107$ \\
  & (IRR\,=\,1.157) & (IRR\,=\,1.113) \\[3pt]
\textbf{Prestige\,$\times$\,ai\_z}
  & $\mathbf{-0.003}$ & $\mathbf{-0.007}$ \\
  & (IRR\,=\,\textbf{0.997}) & (IRR\,=\,\textbf{0.993}) \\[3pt]
\midrule
Controls, Year FEs & Yes & Yes \\
AIC & --- & --- \\
$N$  & 843 & 843 \\
\bottomrule
\end{tabular}
\begin{tablenotes}
\footnotesize
\item \textit{Note.} Coefficients are log-incidence-rate ratios;
  exponentiated IRRs in parentheses.
  Controls: $\log(\text{Age})$, $\log(\text{Abs.\ length})$, year FEs.
  The NB specification shows a significant prestige main effect for
  downloads ($p = .014$) that is not observed in the OLS log(+1)
  specification ($p = .104$), suggesting the prestige effect may be
  non-linear in download scale.
  Crucially, the Prestige\,$\times$\,AI interaction is near zero and
  non-significant in both NB models, consistent with the main OLS results.
  $^{*}p<.05$;\ $^{**}p<.01$.
\end{tablenotes}
\end{threeparttable}
\end{table}

%% ============================================================
\subsection*{B6. Topic Fixed-Effect Sensitivity}
\addcontentsline{toc}{subsection}{B6. Topic Fixed-Effect Sensitivity}

To assess whether the main results are driven by topic-level
heterogeneity, we add binary topic dummies to the main specification.
Topics were assigned using keyword matching on title, abstract, and
keywords: AI/Digital (88.4\%), Supply Chain (2.5\%),
ESG/Green (1.8\%), Marketing (0.8\%), Innovation (3.1\%),
Governance (0.2\%), Other (3.2\%).
Because 88\% of articles fall in a single category (AI/Digital),
the topic dummies explain little additional variance.
Results are nearly identical to the baseline:
$\hat\beta_3 = -0.045$ ($p = .460$) for downloads and
$\hat\beta_3 = +0.003$ ($p = .952$) for citations,
with prestige remaining non-significant ($p = .107$ and $p = .523$).
The focal results are not driven by topic composition.

%% ============================================================
\subsection*{B7. Marketing and Management Sub-Sample}
\addcontentsline{toc}{subsection}{B7. Marketing and Management Sub-Sample}

To address concerns about sample breadth, we restrict the analysis
to articles from CNKI journals whose names contain at least one of
the following management-related terms:
管理 (management), 营销 (marketing), 市场 (market/marketing),
商学 (business), 工商 (business administration),
经济管理 (economic management), 商业 (business),
战略 (strategy), 组织 (organization), 企业 (enterprise/firm).
This yields a sub-sample of $n = 239$ articles (28.4\% of the full sample)
from journals including 管理世界, 南开管理评论, 管理学报,
科研管理, 经济管理, 外国经济与管理, and others.

\begin{table}[htbp]
\centering
\caption{Marketing/Management Sub-Sample ($n = 239$, OLS, HC3 SEs)}
\label{tab:mktgmgmt}
\begin{threeparttable}
\small
\begin{tabular}{lcc}
\toprule
 & (1) Downloads & (2) Citations \\
\midrule
Prestige\_High
  & $-0.009$ & $0.038$ \\
  & (0.152)  & (0.116) \\[4pt]
ai\_z
  & $0.254^{*}$ & $-0.046$ \\
  & (0.117)     & (0.076) \\[4pt]
\textbf{Prestige\,$\times$\,ai\_z}
  & $\mathbf{-0.235}$ & $\mathbf{+0.065}$ \\
  & \textbf{(0.145)}  & \textbf{(0.112)}  \\[4pt]
\midrule
Controls, Year FEs & Yes & Yes \\
$R^2$  & .421 & .702 \\
$N$    & 239  & 239  \\
\bottomrule
\end{tabular}
\begin{tablenotes}
\footnotesize
\item \textit{Note.} OLS with HC3 robust SEs.
  The AI-style main effect on downloads is positive and significant
  ($\hat\beta = 0.254$, $p = .029$), replicating the main finding
  of the full sample.
  The Prestige\,$\times$\,AI interaction is larger in magnitude
  ($\hat\beta_3 = -0.235$) than in the full sample ($-0.044$)
  but remains non-significant ($p = .106$), consistent with
  low statistical power at $n = 239$.
  The prestige main effect is not significant in either outcome.
  $^{*}p<.05$.
\end{tablenotes}
\end{threeparttable}
\end{table}

%% ============================================================
\section*{Appendix C: Comparison with Jo and Li (2026)}
\addcontentsline{toc}{section}{Appendix C: Comparison with Jo and Li (2026)}
%% ============================================================

Table~\ref{tab:compare} places the present study's focal estimates
alongside the corresponding estimates from Jo and Li (2026) for
direct comparison.
The two studies differ in data source (SSRN preprints vs.\ CNKI
peer-reviewed articles), language (English vs.\ Chinese),
AI-style measurement (LLM-based evaluation vs.\ dictionary marker count),
prestige measurement (QS rankings vs.\ C9/985/211 tier), and
outcome metrics (SSRN abstract views vs.\ CNKI downloads/citations).
These design differences preclude direct statistical comparison of
point estimates; the table is intended to summarize the direction and
rough magnitude of replicated and non-replicated effects.

\begin{table}[htbp]
\centering
\caption{Side-by-Side Comparison with Jo and Li (2026)}
\label{tab:compare}
\begin{threeparttable}
\small
\begin{tabular}{p{4.5cm}cc}
\toprule
& Jo \& Li (2026) & Present study \\
& SSRN, English & CNKI, Chinese \\
\midrule
\textit{Sample} & & \\
\quad Context & Preprints, no gatekeeping & Peer-reviewed journals \\
\quad$N$ & [not disclosed] & 843 \\
\quad Period & [not disclosed] & 2023--2026 \\[4pt]
\textit{AI-style $\to$ Diffusion} & & \\
\quad Direction & Positive & Positive \\
\quad Significant (downloads) & Yes & Yes ($p < .001$) \\
\quad Significant (citations) & Yes & Marginal ($p = .022$--$.178$) \\[4pt]
\textit{Prestige $\to$ Diffusion} & & \\
\quad Direction & Positive & Positive \\
\quad Significant & Yes & No ($p = .104$--$.606$) \\[4pt]
\textit{Prestige\,$\times$\,AI interaction} & & \\
\quad Direction & Negative (attenuating) & Negative or near zero \\
\quad Significant & Yes & No (all $p > .28$) \\
\quad 95\% CI (download model) & [not yet available] & $[-0.164, +0.077]$ \\[4pt]
\textit{Boundary condition} & & \\
\quad Venue type & Open preprint & Peer-reviewed \\
\quad Prestige premium & Operative & Absent (null in OLS) \\
\quad AI moderation & Significant & Moot (no premium to moderate) \\
\bottomrule
\end{tabular}
\begin{tablenotes}
\footnotesize
\item \textit{Note.} Jo and Li (2026) coefficient details will be inserted
  after the paper is formally published with volume/issue/page numbers.
  ``Moot'' in the present study reflects that the moderation mechanism
  requires an operative prestige premium as a precondition.
\end{tablenotes}
\end{threeparttable}
\end{table}

\end{document}
